\chapter{Introduction}

\section{Background}

Forest fires are one of the most serious natural hazards affecting Türkiye.
Every year, thousands of hectares of forest are damaged, threatening human
lives, wildlife, critical infrastructure, and the national economy~\cite{colak2020}.
The combination of hotter summers, prolonged drought periods, and human
activities has increased both the frequency and severity of wildfires~\cite{goltas2024,ekberzade2025}.

Traditional wildfire monitoring in Türkiye still relies heavily on fixed
lookout towers, ground patrols, and reports from local communities. Although
these methods are useful, they suffer from several limitations: they are
labor-intensive, dependent on human vigilance, and often detect fires only
after they have grown to a dangerous size~\cite{bugaric2025}. At the same
time, recent advances in unmanned aerial vehicles (UAVs), low-cost sensors,
and artificial intelligence (AI) provide a new opportunity to move from
reactive firefighting to proactive, risk-driven monitoring~\cite{mohapatra2022}.

This project focuses on FireWatch: A Drone-Ready Wildfire Risk Prediction and
Fire/Smoke Detection System for the Karabük region. The final system combines
daily wildfire risk prediction from 2012--2025 seasonal Karabük records, live
fire/smoke detection using a custom YOLOv8 detector, persistent alert storage,
and dashboard-based monitoring. The system also includes a drone-ready layer
for mock, Tello, and demo support, but it does not claim to be a fully
autonomous production drone deployment.

\section{Problem Statement}

Despite the increasing impact of wildfires in Türkiye, current monitoring and
early-warning practices still show important gaps:

\begin{itemize}
  \item Fire risk is not always quantified in a continuous and operational form
  that can support daily monitoring decisions for the target region~\cite{colak2020}.
  \item Drone and camera systems, when used, are usually operated manually
  rather than being prioritized or recommended based on risk levels~\cite{rjoub2023}.
  \item Many existing camera systems only provide live video without
  intelligent fire/smoke detection, which means they still depend on human
  operators~\cite{mohapatra2022}.
  \item Historical data on risk, detections, and responses are not fully
  integrated into a single platform that can support post-event analysis and
  future planning.
\end{itemize}

As a result, fires may be detected late, high-risk periods may not receive
enough attention, and decision-makers may lack an integrated view of risk,
detection, and response. In the final FireWatch system, this gap is addressed
by combining risk prediction, live fire/smoke detection, persistent alerts, and
a dashboard view.

Prediction alone cannot confirm whether fire or smoke is already present, while
visual detection alone does not provide the daily risk context needed for
preparedness and prioritization. Operators therefore need an integrated system
that connects risk estimates, visual evidence, alert history, and system status
in one workflow.

\section{Importance of the Study}

Developing an intelligent wildfire monitoring system has clear real-world
value. Earlier and more reliable detection of fires can reduce burned area,
economic losses, and ecological damage~\cite{mohapatra2022}. A risk-driven
monitoring approach can help forestry authorities and operators focus limited
resources on the most vulnerable periods and locations instead of monitoring
large areas uniformly.

The drone-ready layer adds another engineering dimension by preparing the
system for future patrol support. In the current implementation, this layer is
limited to mock, Tello, and demo support and should be understood as a
prototype direction rather than a field-tested autonomous drone deployment.

\section{Research Question}

The main research question of this thesis is:

\begin{quote}
How can a drone-ready wildfire monitoring system combine daily wildfire risk
prediction, live fire/smoke detection, persistent alert management, and
dashboard-based decision support for the Karabük region?
\end{quote}

\section{Aim and Objectives}

The aim of this project is to design and implement FireWatch as a drone-ready
wildfire risk prediction and fire/smoke detection prototype for the Karabük
region.

To achieve this aim, the project sets the following specific objectives:

\begin{enumerate}
  \item Predict FWI and high risk days using machine learning instead of traditional calculation methods.
  \item Flag high-risk days early using a clear threshold based on historical fire records.
  \item Detect and confirm fire/smoke in real time using a custom YOLOv8 model.
  \item Integrate prediction, detection, alerts, and dashboard monitoring into one unified system.
\end{enumerate}

Each objective directly addresses limitations identified in the problem
statement and can be evaluated using measurable indicators such as prediction
performance, detection performance, response time, persistence behavior,
testing results, and dashboard usability.

\section{Scope of the Project}

This project is scoped as a pilot implementation for \textbf{Karabük Province}
(Türkiye) and its surrounding forested areas. Karabük is located in the Western
Black Sea region and is characterized by extensive forest coverage and mixed
terrain, where seasonal heat, dry winds, and human activity near roads and
settlements can elevate wildfire risk.

The project focuses on the period from 2012 to 2025, which includes 14 fire seasons
and provides a rich dataset for training and testing the prediction model.
\section{Main Contributions}

The main contributions of this thesis are summarized as follows:

\begin{itemize}
  \item An end-to-end FireWatch prototype that combines wildfire risk
  prediction, fire/smoke detection, alert storage, and dashboard monitoring.
  \item A 34-feature prediction pipeline for regional daily wildfire risk
  estimation.
  \item A two-stage prediction model using a Stage 1
  HistGradientBoostingRegressor and a Stage 2 RandomForestClassifier safety
  classifier.
  \item A custom YOLOv8 fire/smoke detector using \texttt{best.pt} as the
  active detection model.
  \item A dashboard that includes Live Monitoring, Risk Decision, Risk Map,
  Detection Alerts, System and Model, and System Flow views.
\end{itemize}

\section{Thesis Organization}

Chapter 2 reviews related studies on wildfire prediction and fire/smoke
detection, including the limitations of reviewed studies and a comparative
summary. Chapter 3 explains the methodology, including data preparation,
feature engineering, prediction modeling, and visual detection methodology.
Chapter 4 presents the system design and architecture, including the backend,
frontend, database, scheduler, and drone-ready layer.

Chapter 5 presents the results and discussion, including prediction results,
YOLOv8 detection evidence, dashboard behaviour, validation, limitations, and
interpretation. Chapter 6 concludes the thesis and describes future work.